\section{The formalised abstract interface}\label{app:lean}

The relevant source is \texttt{NoFiniteStencilMinorant.lean} in
\texttt{Q3/Proofs/RouteB/}, in repository \texttt{Malaeu/chen\_q3},
at the manuscript source revision \texttt{0f9c7cbc81bb}.
Its Git blob begins \texttt{e6fd9dc2f713}.
The following five exports are printed in the file's axiom audit:
\begin{itemize}[leftmargin=2em]
\item \texttt{independence\_of\_translates};
\item \texttt{independence\_of\_translates\_apply};
\item \texttt{stencil\_energy\_limit\_eq\_zero};
\item \texttt{no\_positive\_finite\_stencil\_minorant};
\item \texttt{no\_positive\_finite\_stencil\_minorant\_hypotheses\_satisfiable}.
\end{itemize}
Five is the number of these audited exports, not the number of all
helper declarations. The proof includes real and complex analysis,
Fourier integrals, Fatou's lemma and almost-everywhere reasoning as
well as finite Vandermonde algebra.

In the principal theorem, $f:\R\to\R$ is continuous, integrable and
strictly positive; $Q:(\R\to\C)\to\R$ is arbitrary; the cutoffs are
measurable and equal to $1$ on $[-R,R]$; the real coefficient stencil
is nonzero and its shifts are distinct. The weight $W$ takes values
in the extended nonnegative reals. With
$r_q(x)=f(x-q)/f(x)$, the hypotheses are:
\[
 \text{(H1)}\quad Q(f\chi_Rr_q)\longrightarrow0
       \quad\text{for every rational }q,
\]
\[
 \text{(H2)}\quad
 \int W|S(\chi_Rr_q)|^2
       \le\max\{Q(f\chi_Rr_q),0\}
       \quad(q\in\mathbb Q,\ R\ge1),
\]
where the integral and inequality in (H2) are represented in
$[0,\infty]$. In Lean the right-hand side is
\texttt{ENNReal.ofReal}. The conclusion is $W=0$ almost everywhere.
The real-valued minorant inequality in \Cref{thm:nominorant} implies
this formal hypothesis, but is not literally the same expression.
For the concrete Weil application (H1) is supplied by the paper
proof of \Cref{lem:budget}, not by this Lean file.

The independence result allows complex coefficients and proves the
finite exponential-polynomial step by sampling and a Vandermonde
matrix. The nonvacuity export uses a Gaussian, the zero functional
and zero weight; it is not a nonzero Weil certificate. The recorded
axiom profiles of the five exports are
\texttt{[propext, Classical.choice, Quot.sound]}.
No claim is made that the explicit formula, the canonical analytic
test, its concrete cutoff budget or the ground-state identity has
been formalised by these exports. Companion matrix and transform
files are outside this paper's formal theorem claim.
